\subsection{The actual CM pullback and its common covariance}

\begin{result}{E.190}{Proved}
{The division quartic gives one comparison in the ordered-pole group}
Let \(F=(1,f)\) be a full GT element with \(\rho(F)=0\), as
in E.181.  Put
\[
 \mathcal B=\widehat{C_2(s)*C_3(r)},\quad
 \tau=sr,\quad x=\tau^{-2},\quad y=sxs,
 \qquad
 G_2=\widehat\Gamma_{1,2},\quad
 G_1=\widehat\Gamma_{1,1}.
\tag{E.190.1}
\]
Use the cusp-normalized source action \(\alpha\) and its
literal pure frames
\[
 \alpha(\tau)=\tau,\qquad
 \alpha(s)=\beta s\beta^{-1},\qquad
 \alpha(r)=ArA^{-1},\qquad A,\beta\in\widehat F(x,y).
\tag{E.190.2}
\]
Let \(\Phi\) be the standard four-strand action on \(G_2\),
with the same parameter \(F\), and let \(q:G_2\to G_1\)
forget the second ordered point.  There is an actual
geometric isomorphism
\[
 \varphi:\mathcal B\xrightarrow{\ \sim\ }
 G_3:=q^{-1}(\langle R_+\rangle),\qquad
 R_+=uw\in G_1,\quad |R_+|=6,
\tag{E.190.3}
\]
whose restriction identifies \(\mathcal B'\) with
\(K=\ker q\), including their full profinite topologies.
It sends the generator \(\tau\) of the cyclic quotient
to \(R_+^\varepsilon\), for one
\(\varepsilon\in\{1,-1\}\).

Lift the actual word \(A(x,y)\) to
\(\widetilde A=A(a^2,b^2)\in G_2\), where \(a,b,c\)
are the consecutive four-strand generators and
\(q(a)=u,\ q(b)=w,\ q(c)=u\).  Then
\[
 \Phi^A=\operatorname{Ad}(\widetilde A^{-1})\Phi
 \quad\hbox{preserves }G_3,
 \qquad
 \Phi^A\varphi=\operatorname{Ad}(h)\varphi\alpha
 \quad\hbox{for one }h\in G_3.
\tag{E.190.4}
\]
In particular the same element \(h\), and not independently
selected conjugators in coefficient quotients, controls all
homological reductions of this equation.
\end{result}

\begin{proof}
\emph{Retaining the actual source subgroup.}
There is a necessary retention step between E.188 and the
CM map.  Let \(\mathcal L_9\) be the ordinary linear
zero-stabilizer of the three-torsion grid, and let
\(\mathcal L_{12}\) be its canonical linear lift obtained
by marking the other three two-torsion points.
Ordinary forgetting identifies both groups with
\(\operatorname{SL}_2(\mathbb Z)\); both have full induced
topology, because their projection to the one-point group
is a retraction on the respective zero-label stabilizer.
Write \(\iota\) for the linear elliptic involution.

The relative kernel of forgetting the three new ordered
marks is the full completion \(\widehat K_{\rm rel}\) of
the fundamental group of
\[
       Y=\operatorname{Conf}_3(T\setminus T[3]).
\tag{E.190.5}
\]
Its full topology and the exact completed forgetting
sequence follow from the iterated actual point-pushing
comparison of E.182.  The fixed locus \(Y^\iota\) consists
of the six isolated assignments of the nonzero
two-torsion points.  A loop in its base component is
therefore trivial.  The one-cyclic fixed-path argument
of E.176, applied exactly as in E.188, gives
\[
                 C_{\widehat K_{\rm rel}}(\iota)=1.
\tag{E.190.6}
\]
This conclusion uses its finite centralizer projections
and compact passage; it is not inferred merely from the
ordinary fixed locus.

We also need this conclusion with the three new labels
unordered.  Apply the same finite theorem to the ordered
space \(Y\), now with its commuting finite actions
\(C_2\times\mathfrak S_3\).  Its actual orbifold path
extension permits permutations of the new labels.
A fixed path starting at the chosen assignment must
end at that same point of the six-point fixed locus.
A nonidentity label permutation has a different
endpoint; the identity endpoint has no nonconstant
fixed loop.  The finite centralizer recognition
therefore gives
\[
       C_{\widehat K_{\rm rel}^{\,\rm unord}}(\iota)=1
\tag{E.190.6a}
\]
for the relative kernel with these labels unordered.
This uses the ordered poly-free cover and the full
finite group of symmetries, not a poly-free assertion
about an unordered braid group.

The corrected twelve-point action fixes \(\iota\) and
has the corrected nine-point action as its projection.
For \(g\in\widehat{\mathcal L}_{12}\), its image and
the canonical linear lift of its nine-point image
therefore centralize \(\iota\) and have the same
nine-point projection.  Their quotient lies in the
unordered relative kernel in E.190.6a, so they are equal.
In particular this argument determines their common
permutation of the new labels; it does not assume it
from the old nine-label permutation alone.
Thus the action retains the actual canonical linear
lift, and on it has precisely the original one-point
outer action of E.184.

\emph{The map after the quotient and forgetting.}
Divide the twelve-point surface by \(\iota\), then
forget its three nonzero fixed branch labels.
The remaining marks are the origin and the four
images of the opposite nonzero three-torsion pairs.
E.189 proves that this actual five-point group has
the original parameter \(F\), with the origin retained.
The source \(\mathcal L_{12}\) fixes the origin and
permutes the other four retained marks.  Its central
element \(-I\) becomes the deck involution and dies
on the quotient.  Hence its induced map factors
through the actual rigidified one-point modular group
\(\mathcal B\).

We specify why the comparison also holds on the
mixed target, rather than only on its pure subgroup.
Let \(D_5\) denote the sphere group with its origin
fixed and its four other marks unordered.  Its pure
subgroup is \(\Gamma_{0,5}\), and
\[
                 D_5/\Gamma_{0,5}=\mathfrak S_4.
\tag{E.190.7}
\]
The ordinary forgetful map from the origin-stabilizer
in the colored eight-point group to \(D_5\) is onto.
Its kernel is the ordinary group of motions and
permutations of the three forgotten branch marks.
The action on its pure part has the standard
forgetting covariance by E.189.  After completion,
any image of the whole kernel centralizes the
pure target: its commutators with the pure source
belong to the pure kernel.
The centralizer of \(\widehat\Gamma_{0,5}\) in
\(\widehat D_5\) is trivial.  Indeed its permutation
is trivial by the faithful geometric \(\mathfrak S_5\)
factor of the HMM theorem used in E.189; it then
lies in the center of \(\widehat\Gamma_{0,5}\).
That center is zero: forget to
\(\widehat\Gamma_{0,4}=\widehat F_2\), then use the
center-free free kernel \(\widehat F_3\).
Consequently the completed forgetful kernel is
preserved and the mixed action descends.

The same argument proves uniqueness of the mixed
extension of a specified pure action.  After one
inner correction the restrictions of two extensions
to the pure subgroup agree exactly; their values on
any mixed element differ by an element centralizing
that pure subgroup, hence agree.  Thus the action
on \(\widehat D_5\) is the standard action with
parameter \(F\), up to one inner change.
The geometric permutation ambiguity of E.189 is
supported on the three forgotten labels and has
already disappeared.  E.183 identifies this mixed
five-point group with \(G_2\) through the double
cover with its two poles ordered.  All these maps
are the indicated ordinary maps before completion.

\emph{The monic division quartic and its poles.}
For the Weierstrass curve
\(y^2=x^3+ax+b\), set \(D=4a^3+27b^2\ne0\).
The roots of its three-division polynomial
\(3x^4+6ax^2+12bx-a^2\) are precisely the four
pair labels.  Taking the double cover branched at
those four marks gives
\[
 \begin{gathered}
 t^2=x^4+2ax^2+4bx-a^2/3,\qquad
 X=(t+x^2+a)/2,\\
 Z=X-a/3,\qquad Y=Xx+b/2,\\
 Y^2=Z^3+D/108,\qquad P=(-a/3,-b/2),\\
 x=\frac{Y-b/2}{Z+a/3},\qquad
 t=2(Z+a/3)-x^2-a.
 \end{gathered}
\tag{E.190.8}
\]
Substitution proves the two inverse identities in
the function fields; smooth projective completion
extends them across the omitted points.
The pole with \(t/x^2=+1\) is the origin \(O\);
the other pole is \(P\).
The degree-two function \(x\) has pole divisor
\(O+P\), so its deck involution is \(Q\mapsto P-Q\).
At \(b=0\), \(P\) is a nonzero point of order two;
at \(a=0\), its horizontal tangent has intersection
multiplicity three, so it is a nonzero point of
order three.

These are ordered poles intrinsically.  Under
\(x'=cx+d\) the monic square root is changed by
\(t'=c^2t\), retaining the sign of \(t/x^2\).
Permuting the four roots also preserves the monic
polynomial and each of those signs.
Thus this construction lands in the ordered group
\(G_2\), with the marking used in E.183.

Original Weierstrass scaling by \(u\) induces
\((Z,Y)\mapsto(u^4Z,u^6Y)\).  After the original
central involution is rigidified, the residual
variable \(v=u^2\) has weights \(2,3\) on \(a,b\),
and \(D\) has weight six.  On the slice \(D=108\)
put \(a=-3Z,\ b=-2Y\).  This gives the actual
stack isomorphism
\[
 \left[\{4a^3+27b^2\ne0\}/\mathbb G_m{}_{(2,3)}\right]
       \simeq
 \left[(E_0\setminus\{O\})/\mu_6\right],
 \qquad E_0:\ Y^2=Z^3+1 .
\tag{E.190.9}
\]
Indeed choosing a sixth root of \(D/108\) normalizes
to that slice, and two choices differ exactly by
\(\mu_6\), acting with the displayed weights.
The left orbifold is the rigidified one-point
elliptic moduli stack, with orbifold orders two
and three and a cusp, hence fundamental group
\(C_2*C_3\).  The right side is the actual
universal punctured elliptic curve over its
order-six CM stabilizer.  Its pullback group is
the ordinary preimage of \(\langle uw\rangle\)
under \(\Gamma_{1,2}\to\Gamma_{1,1}\).

This proves the discrete version of E.190.3.
Its kernel over \(C_6\) is the actual punctured
torus group \(F_2\), and on the source it is the
index-six commutator subgroup of \(C_2*C_3\).
The source kernel has full topology by finite
index.  The target kernel has full topology by
E.182; right exactness of completion and that
topology give the exact completed Birman
sequence.  The completed map is an isomorphism
on its kernel and on its finite quotient,
therefore on the whole pullback.
In particular \(G_3\) has not been treated as
an open subgroup of \(G_2\).
The quotient isomorphism sends its generator
to \(R_+^\varepsilon\), with
\(\varepsilon=\pm1\); its sign will not be
discarded.

The retained linear subgroup, quotient, forgetting
and ordered-pole constructions therefore give one
outer-commutative square
\(\Phi\varphi\simeq\varphi\alpha\).
Each inner source change maps through \(\varphi\)
and can be absorbed into its one target conjugator.
There has been no choice of separate coefficient
or pairwise final conjugators.

\emph{The CM frame and location of the conjugator.}
The standard action on \(G_1\), projected modulo
its central involution, is the source action
E.190.2 in the literal modular marking of
E.180.  Thus
\[
       \Phi(R_+)=q(\widetilde A)R_+
                       q(\widetilde A)^{-1}
\tag{E.190.10}
\]
up to that central involution.  The latter
ambiguity is absent: \(G_1^{\rm ab}=\mathbb Z/12\),
the multiplier-one action is identity on this
abelianization, and multiplying by the central
involution changes its class by six.
This proves E.190.10 exactly, so \(\Phi^A\)
fixes \(R_+\) in the quotient and preserves \(G_3\).
The common covariance already proved gives
E.190.4 with an element initially in \(G_2\).
Both maps induce identity on the CM cyclic
quotient, so \(q(h)\) centralizes \(R_+\).
Modulo the central involution this is the
centralizer of the order-three free factor
in \(\widehat{C_2*C_3}\), which is precisely
that factor by
\citet[Theorems A--B and B', pp.155--160]
{SymHerfortRibes1985}.
Its full preimage in \(G_1\) is
\(\langle R_+\rangle\).
Hence \(h\in G_3\), proving all assertions.
\end{proof}
